% Marshall Bruner's resume
%
% HOW TO EDIT
% - Each entry is one macro call: \job, \project, \projectsingle, \school, \skill.
% - Bullets are a comma-separated list. Wrap EVERY bullet in its own braces
%   so commas inside the bullet text don't split it:
%     { {First bullet, with a comma}, {Second bullet} }
% - Bullets may wrap onto several lines; spacing is handled by the macros.
%   Do not add manual \vspace between entries or sections.
% - Entry headers never get orphaned at a page bottom (needspace).
% - \boldblue{...} highlights text; remove the wrapper to un-highlight.
% - Compile: pdflatex resume.tex

\documentclass[11pt,letterpaper]{article}

\usepackage[english]{babel}
\usepackage[utf8]{inputenc}
\usepackage{fontawesome5}
\usepackage{plex-mono}   % IBM Plex Mono for the name (\plexmono)
\usepackage[hidelinks]{hyperref}
\usepackage{tabularx}
\usepackage{listofitems} % comma-separated bullet lists
\usepackage{xcolor}
\usepackage{enumitem}    % tight bullet lists
\usepackage{needspace}   % keep entry headers with their bullets
\usepackage{titlesec}    % section spacing and rule

\usepackage{geometry}
\geometry{
  letterpaper,
  left=6mm,
  right=6mm,
  top=10mm,
  bottom=10mm
}

\pagestyle{empty}
\setlength{\parindent}{0pt}

% ---------- Spacing knobs ----------
\newcommand{\entrygap}{3.5mm}      % space after each entry (separates jobs)
\newcommand{\headtobullets}{-2.5mm} % pull bullets up under the entry header
\setlist[itemize]{leftmargin=1.6em, nosep, topsep=0pt, itemsep=1pt}

% ---------- Section formatting ----------
\titleformat{\section}{\large\bfseries}{}{0pt}{}[\vspace{-7pt}\rule{\textwidth}{0.5pt}]
\titlespacing*{\section}{0pt}{5pt}{3pt}

% Highlight command. Pick one definition:
\DeclareRobustCommand{\boldblue}[1]{#1}                               % off (plain text)
% \DeclareRobustCommand{\boldblue}[1]{\textbf{#1}}                    % bold only
% \DeclareRobustCommand{\boldblue}[1]{{\textcolor{blue}{\textbf{#1}}}} % bold + blue

% ---------- Header ----------
% Arguments: name, phone, email, linkedin handle, github handle, youtube handle, website
\newcommand{\namefull}[7]{
  \begin{center}
    {\Huge \plexmono #1} \\ \vspace{3mm}
    \small \raisebox{-0.1\height}\faPhone \ \underline{#2} ~
    \href{mailto:#3}{\raisebox{-0.1\height} \faEnvelope \  \underline{#3}} ~
    \href{https://linkedin.com/in/#4}{\raisebox{-0.1\height} \faLinkedin \  \underline{#4}} \\ \vspace{0.5mm}
    \href{https://github.com/#5}{\raisebox{-0.1\height} \faGithub \  \underline{#5}} ~
    \href{https://youtube.com/#6}{\raisebox{-0.1\height} \faYoutube \  \underline{#6}} ~
    \href{https://#7}{\raisebox{-0.1\height} \faGlobe \  \underline{#7}}
  \end{center}\par
}

% Header without youtube / website (kept for older branches)
% Arguments: name, phone, email, linkedin handle, github handle, (unused)
\newcommand{\namenoloc}[6]{
  \begin{center}
    {\Huge \plexmono #1} \\ \vspace{4mm}
    \small \raisebox{-0.1\height}\faPhone \ \underline{#2} ~
    \href{mailto:#3}{\raisebox{-0.1\height} \faEnvelope \  \underline{#3}} ~
    \href{https://linkedin.com/in/#4}{\raisebox{-0.1\height} \faLinkedin \  \underline{#4}} ~
    \href{https://github.com/#5}{\raisebox{-0.1\height} \faGithub \  \underline{#5}}
  \end{center}\par
}

% ---------- Entry macros ----------
% Shared bullet list: argument is the comma-separated list
\newcommand{\bullets}[1]{%
  \readlist*\bulletslist{#1}%
  \vspace{\headtobullets}
  \begin{itemize}
    \foreachitem\bullet\in\bulletslist{\item \bullet}
  \end{itemize}
  \vspace{\entrygap}
}

% Job
% Arguments: company, start date, end date, position, location, bullets
\newcommand{\job}[6]{
  \Needspace*{5\baselineskip}
  \begin{tabularx}{\textwidth}{
    p{14.5cm}
    >{\raggedleft\arraybackslash}X
    }
    {\large\textbf{#1}} & \textbf{#2} \textemdash\ \textbf{#3} \\
    \textit{\underline{#4}} & \textit{#5} \\
  \end{tabularx}
  \bullets{#6}
}

% Project with a date range
% Arguments: title, start date, end date, bullets, link (may be empty)
\newcommand{\project}[5]{
  \Needspace*{4\baselineskip}
  \begin{tabularx}{\textwidth}{
    p{14.5cm}
    >{\raggedleft\arraybackslash}X
    }
    \href{#5}{\underline{{\large\textbf{#1}}}} & \textbf{#2} \textemdash\ \textbf{#3} \\
  \end{tabularx}
  \bullets{#4}
}

% Project with a single date
% Arguments: title, date, bullets, link (may be empty)
\newcommand{\projectsingle}[4]{
  \Needspace*{3\baselineskip}
  \begin{tabularx}{\textwidth}{
    p{14.5cm}
    >{\raggedleft\arraybackslash}X
    }
    \href{#4}{\underline{{\large\textbf{#1}}}} & \textbf{#2} \\
  \end{tabularx}
  \bullets{#3}
}

% School
% Arguments: school, degree, GPA, start date, end date, note (thesis etc.), location
\newcommand{\school}[7]{
  \Needspace*{4\baselineskip}
  \begin{tabularx}{\textwidth}{
    p{14.5cm}
    >{\raggedleft\arraybackslash}X
  }
    {\large\textbf{#1}} & \textbf{#4 \textemdash\ #5} \\
    \textit{#2} \textemdash\ \textbf{\textit{GPA: #3}} & \textit{#7} \\
  \end{tabularx}
  \if\relax\detokenize{#6}\relax\else
    \par\vspace{-1mm}\hspace{\tabcolsep}\textit{#6}\par
  \fi
  \vspace{\entrygap}
}

% Skill
% Arguments: group, skills
\newcommand{\skill}[2]{
  \begin{tabularx}{\textwidth}{>{\raggedright\arraybackslash}X}
    \textbf{#1:} #2
  \end{tabularx}
}

% =====================================================================
\begin{document}
\namefull{Marshall Bruner}
{+1 719-553-9497}
{brunerm99@gmail.com}
{marshall-bruner}
{brunerm99}
{@MarshallBrunerRF}
{marshallbruner.com}

\vspace{-3mm}
\begin{center}
  \textit{RF hardware, signal processing, and weather radar deployment \textbar\ Teaching 50k+ engineers on YouTube}
\end{center}

\section*{EXPERIENCE}
\job{Rainmaker Technology Corporation}
{Aug 2025}
{Oct 2026}
{Head Radar Systems Engineer}
{Boulder, CO}
{
  {\boldblue{Led radar engineering for a cloud-seeding company}: set S-, X-, and K-band sensor strategy and owned radar hardware, FPGA / embedded software, signal processing, and deployment tooling},
  {\boldblue{Designed a 24 GHz FMCW cloud-profiling radar transceiver} (8-layer Rogers PCB, 450+ components, KiCad) from architecture through fab and bring-up, including a bistatic redesign for 15+ dB sensitivity gain and GaN PA bias sequencing},
  {Owned its software stack: SystemVerilog gateware on Zynq UltraScale+ (PLL control, LVDS ADC capture, DMA), C++ control server, Python capture and Doppler processing pipeline, and SCPI-automated bench characterization},
  {\boldblue{Operated and deployed a 25 kW X-band dual-pol scanning weather radar}: scan-control and monitoring service with Slack / SNMP alerting, remote networking, and a mountain-top installation},
  {Architected a solid-state pulse-compression X-band dual-pol successor radar (transceiver design, PA trade study, sensitivity budget, business case) and built terrain-aware coverage-planning and radar performance-modeling tools used to site radars in 12 program regions}
}
\job{CSU-CHILL National Radar Facility}
{Jan 2025}
{Present}
{Radar Systems Engineer (part-time)}
{Greeley, CO}
{
  {\boldblue{Built the facility's multi-radar web display} (Flask, OpenLayers, PostgreSQL, Docker) serving real-time and historical PPI / RHI views for 8 radars and 14 polarimetric products, with client-side GPU rendering and a multi-layer cache},
  {\boldblue{Wrote the antenna positioner controller for a solid-state C-band radar} in Go: reverse-implemented the pedestal's binary UDP protocol, PPI / RHI scan planner and queue, operator web UI, and CLI},
  {Maintained an X-band dual-pol signal processing chain in C (time series to moments to CF-Radial NetCDF to polarimetric products) and authored its operational volume coverage patterns},
  {Supported relocation and re-commissioning of a dual-pol S-band weather radar at Reagan Test Site, Kwajalein Atoll: bring-up procedure, calibration verification, and a live monitoring dashboard}
}
\job{BAE Systems (\boldblue{DoD Clearance})}
{Jul 2024}
{Jul 2025}
{RF Engineer}
{Westminster, CO}
{
  {\boldblue{System design, schematic capture, and layout of a 2D active phased array weather radar}},
  {Large-scale phased array systems analysis in Python and Matlab},
  {Technical consultant on weather radar proposals}
}
\job{Ball Aerospace (\boldblue{DoD Clearance})}
{Jun 2021}
{Jul 2023}
{RF Engineer}
{Westminster, CO}
{
  {\boldblue{Designed and debugged multiple mission-critical RF boards under strict time constraints}},
  {Simulated, designed, and tested an RF front-end PCB at extreme temperatures},
  {\boldblue{Schematic capture and layout of RF T/R systems} (RF, digital, power), antenna range test equipment, and component test coupons},
  {Built an embedded webpage used to control and monitor phased array electronics; high-sensitivity noise and high-power testing}
}
\job{Colorado State University}
{Jan 2021}
{May 2022}
{Graduate Research Assistant | Thesis focused on Radar}
{Fort Collins, CO}
{
  {\boldblue{Design, schematic, and layout} of a modular X-band FMCW drone radar},
  {\boldblue{Signal processing and data visualization} for Analog Devices' FMCW Phaser radar board; presented by ADI at IMS 2022 and IEEE-APS 2022},
  {Presented the work at the IEEE International Symposium on Phased Array Systems and Technology}
}
\job{CSU-CHILL Radar Lab}
{May 2020}
{May 2021}
{Lab Assistant}
{Greeley, CO}
{
  {Built a server to host and display real-time and historical radar data using PostgreSQL, Flask, and Python},
  {Performed characterization and measurements of radar equipment}
}

\section*{PROJECTS}
\project{Marshall Bruner RF \textemdash\ Radar / RF Education YouTube Channel}
{Aug 2023}
{Present}
{
  {\boldblue{53k subscribers and 1.4M+ views} across 17 long-form animated videos on FMCW radar, CFAR, phased arrays, pulse compression, I/Q sampling, and dual-pol weather radar; top video 320k+ views},
  {Write, animate (Manim), narrate, and edit every video; built a reusable radar animation library (WSR-88D model, PPI displays, beam geometry, NEXRAD ingest) and contributed features to the ManimGX GPU engine},
  {Ship companion interactive notebooks (marimo to WASM) and a Python / Three.js 3D weather-radar animation studio with DEM terrain and 4/3-Earth beam refraction}
}
{https://www.youtube.com/@MarshallBrunerRF}
\projectsingle{FMCW Radar Bootcamp \textemdash\ IEEE AESS RadarConf}
{Mar 2023}
{{Developed \boldblue{signal processing software and a presentation for an FMCW radar} for the 2023 AESS Radar Bootcamp}}
{https://radar2023.ieee-radarconf.org/radar-boot-camp}

\section*{EDUCATION}
\school{Colorado State University}
{M.S. Electrical Engineering}
{3.9}
{Jan 2021}
{May 2022}
{Thesis - Design, Deployment, and Cost Considerations for DARMA; A Low-Cost and Lightweight FMCW Radar}
{Fort Collins, CO}
\school{Colorado State University}
{B.S. Electrical Engineering}
{3.47}
{Aug 2017}
{May 2021}
{}
{Fort Collins, CO}

\section*{SKILLS \& ABILITIES}
\skill{Radar / RF}
{FMCW and pulsed radar design, pulse compression, dual-polarization polarimetry, phased arrays, link budgets, Doppler / spectral moment estimation, CFAR, receiver noise and sensitivity analysis, weather radar operations and deployment}
\skill{Languages}
{Python (NumPy, SciPy, Pandas, PyTorch, FastAPI, Flask), Rust, Go, C / C++, SystemVerilog, Bash, SQL, Matlab, TypeScript}
\skill{Design / Simulation Tools}
{KiCad, AWR Microwave Office, HFSS, ADS, LTSpice, Cadence, QucsStudio, Vivado}
\skill{Lab / Infrastructure}
{VNA, spectrum analyzer, oscilloscope, noise figure and high-power test; Git, Linux, Docker, PostgreSQL, Tailscale, \LaTeX, Manim}
\skill{Communication}
{Technical presentations and design reviews (PDR / CDR), customer pre- and post-sales support, educational content reaching 50k+ engineers}

\end{document}
